Partial differential equations (PDEs) are central to scientific computing and engineering, but high-fidelity numerical simulations are often computationally expensive. Neural operators have recently emerged as a promising approach for surrogate modeling by learning mappings between function spaces rather than finite-dimensional vectors.

A key challenge for such models is discretization invariance: the ability to produce consistent predictions when the same underlying problem is represented on different discretizations. This thesis studies this property for Fourier- and convolution-based neural operators on uniform Cartesian grids, and evaluates a geometry-aware Fourier neural operator on non-uniform point-cloud domains.

On uniform grids, the Fourier Neural Operator (FNO) and the Convolutional Neural Operator (CNO) are evaluated on Navier--Stokes and Darcy flow benchmarks through resolution-transfer experiments. The CNO shows substantially flatter loss curves across resolutions, while diagnostic experiments identify fixed-cell spatial padding as the main source of poor FNO resolution transfer in the tested implementation. Resampling inputs to a fixed internal grid largely removes this padding-induced mismatch and makes the FNO loss much more consistent across resolutions.

For non-uniform meshes, the geometry-aware Fourier Neural Operator (Geo-FNO) is evaluated on incompressible fluid flow past a cylinder and on a high-temperature gas-cooled reactor (HTGR) temperature dataset. The results show that Geo-FNO produces accurate predictions on irregular spatial representations, capturing the dominant flow and temperature structures without interpolation onto a uniform Cartesian grid. A full test of mesh-free discretization invariance would require controlled refinement of the same physical samples across multiple point-cloud discretizations.